% TODO make hw stylesheet
\documentclass[10pt]{article}
\usepackage[letterpaper,top=1in]{geometry}
\usepackage[dvipsnames,svgnames]{xcolor}
\usepackage[shortlabels]{enumitem}
\usepackage[framemethod=TikZ]{mdframed}
\usepackage{amsmath,amssymb,amsthm}
\usepackage[colorlinks]{hyperref}
\usepackage{multicol}
\usepackage{tabularx}

\hypersetup{
	colorlinks=true,
	linkcolor=blue,
	filecolor=magenta,
	urlcolor=magenta,
	pdftitle={MAT 534 HW}
}

\usepackage{mathtools}
\usepackage{thmtools}
\usepackage[textsize=footnotesize]{todonotes}
\usepackage{titling}

\reversemarginpar
\mdfdefinestyle{mdbluebox}{roundcorner=10pt,innerbottommargin=9pt,
	linecolor=blue,backgroundcolor=TealBlue!5,}
\declaretheoremstyle[headfont=\sffamily\bfseries\color{MidnightBlue},
mdframed={style=mdbluebox},]{thmbluebox}

\mdfdefinestyle{mdbloobox}{frametitlefont=\bfseries,innerbottommargin=8pt,
	nobreak=true,backgroundcolor=TealBlue!5,linecolor=blue,}
\declaretheoremstyle[headfont=\bfseries\color{MidnightBlue},
mdframed={style=mdbloobox},headpunct={\\[3pt]},postheadspace=0pt,]{thmbloobox}

\mdfdefinestyle{mdredbox}{frametitlefont=\bfseries,innerbottommargin=8pt,
	nobreak=true,backgroundcolor=Salmon!5,linecolor=RawSienna,}
\declaretheoremstyle[headfont=\bfseries\color{RawSienna},
mdframed={style=mdredbox},headpunct={\\[3pt]},postheadspace=0pt,]{thmredbox}

\mdfdefinestyle{mdgreenbox}{linecolor=ForestGreen,backgroundcolor=ForestGreen!5,
	linewidth=2pt,rightline=false,leftline=true,topline=false,bottomline=false,}
\declaretheoremstyle[headfont=\bfseries\sffamily\color{ForestGreen!70!black},
mdframed={style=mdgreenbox},headpunct={ --- },]{thmgreenbox}

\mdfdefinestyle{mdorangebox}{linecolor=RedOrange,backgroundcolor=RedOrange!5,
	linewidth=2pt,rightline=false,leftline=true,topline=false,bottomline=false,}
\declaretheoremstyle[headfont=\bfseries\sffamily\color{RedOrange!70!black},
mdframed={style=mdorangebox},headpunct={ --- },]{thmorangebox}

\mdfdefinestyle{mdblackbox}{linecolor=black,backgroundcolor=RedViolet!5!gray!5,
	linewidth=3pt,rightline=false,leftline=true,topline=false,bottomline=false,}
\declaretheoremstyle[mdframed={style=mdblackbox}]{thmblackbox}

\declaretheoremstyle[spaceabove=6pt,spacebelow=6pt]{basehead}

\declaretheorem[style=basehead,name=Problem]{problem}
\declaretheorem[style=thmredbox,name=Problem,numbered=no]{problem*}
\declaretheorem[style=thmbloobox,name=Setup]{setup} % for config geo
\declaretheorem[style=thmredbox,name=Example,sibling=problem]{example}
\declaretheorem[style=thmredbox,name=$\bigstar$ Problem,sibling=problem]{niceprob}
\declaretheorem[style=thmbluebox,name=Theorem,sibling=problem]{theorem}
\declaretheorem[style=thmbluebox,name=Corollary,sibling=theorem]{corollary}
\declaretheorem[style=thmbluebox,name=Proposition,sibling=theorem]{prop}
\declaretheorem[style=thmbluebox,name=Lemma,numberwithin=problem]{lemma}
\declaretheorem[style=thmbluebox,name=Theorem,numbered=no]{theorem*}
\declaretheorem[style=thmbluebox,name=Lemma,numbered=no]{lemma*}
\declaretheorem[style=thmgreenbox,name=Claim,numberwithin=problem]{claim}
\declaretheorem[style=thmgreenbox,name=Claim,numbered=no]{claim*}
\declaretheorem[style=thmblackbox,name=Remark,numberwithin=problem]{remark}
\declaretheorem[style=thmblackbox,name=Remark,numbered=no]{remark*}
\declaretheorem[style=thmgreenbox,name=Definition,sibling=theorem]{definition}
\declaretheorem[style=thmgreenbox,name=Definition,numbered=no]{definition*}

\providecommand{\ol}{\overline}
\providecommand{\quiv}{\equiv}
\providecommand{\ceil}[1]{\left\lceil #1 \right\rceil}
\providecommand{\floor}[1]{\left\lfloor #1 \right\rfloor}
\newcommand{\dd}{\mathrm{d}}
\providecommand{\ul}{\underline}
\providecommand{\eps}{\varepsilon}
\providecommand{\half}{\frac{1}{2}}
\providecommand{\CC}{\mathbb C}
\providecommand{\FF}{\mathbb F}
\providecommand{\NN}{\mathbb N}
\providecommand{\QQ}{\mathbb Q}
\providecommand{\RR}{\mathbb R}
\providecommand{\ZZ}{\mathbb Z}
\providecommand{\dg}{^\circ}
\providecommand{\ii}{\item}
\providecommand{\setdiff}{\smallsetminus}
\providecommand{\alert}[1]{{\sffamily\textbf{\textcolor{blue}{#1}}}}
\providecommand{\inv}{^{-1}}
\providecommand{\mb}[1]{\mathbf{#1}}

\newenvironment{soln}{\begin{proof}[Solution]}{\end{proof}}

\providecommand{\pow}{\operatorname{Pow}}
\providecommand{\vocab}[1]{{\textbf{\textcolor{ForestGreen}{#1}}}}
\providecommand{\lcm}{\operatorname{lcm}}
\providecommand{\abs}[1]{\left|#1\right|}

\usepackage{mathrsfs}

% place title at top right
\setlength{\droptitle}{-8ex}
\pretitle{\begin{flushright}\Large\bfseries}
\posttitle{\par\end{flushright}}
\preauthor{\begin{flushright}}
\postauthor{\end{flushright}}
\predate{\begin{flushright}}
\postdate{\end{flushright}}

\title{MAT 534 HW03}
\author{\large Aditya Pahuja}
\date{\large Due 09/24}
\begin{document}
\maketitle

\begin{problem}
    Let $G$ be a finite group, $P$ a Sylow $p$-subgroup of $G$.
    If $N$ is a normal subgroup of $G$, show that $P\cap N$
    is a Sylow $p$-subgroup of $N$, and $PN/N$ is a Sylow $p$-subgroup of $G/N$.
\end{problem}

\begin{soln}
    The key fact is that if $H$ is a finite group and $Q\le H$
    is a subgroup with order a power of $p$ such that $(H:Q)$ isn't divisible by $p$,
    then $Q$ is a Sylow $p$-subgroup of $H$ (since its order is therefore
    the largest power of $p$ that divides $\abs H$).

    By HW02 \#1 we have
    \begin{equation*}
	(N : P\cap N) = (PN : P) \mid (G : P).
    \end{equation*}
    Since the order of $P$ is the largest power of $p$ dividing $G$,
    the index of $P$ in $G$ is not a multiple of $p$,
    hence neither is $(N : P\cap N)$. Since $P\cap N$ is a subgroup of $P$
    it has prime power order, so these two facts together imply that
    $P\cap N$ is a Sylow $p$-subgroup of $N$.

    Furthermore, $(PN : N) = (P : P\cap N)$ is a power of $p$,
    so $PN/N$ has order a power of $p$, and
    \begin{equation*}
	(G/N : PN/N) = (G : PN) \mid (G : P)
    \end{equation*}
    is not a multiple of $P$, so $PN/N$ is a Sylow $p$-subgroup of $G/N$.
\end{soln}

\begin{problem}
    Let $G$ be a finite group, $p$ a prime, and $H$ a $p$-subgroup of $G$.
    Show that there is a Sylow $p$-subgroup of $G$ containing $H$.
\end{problem}

\begin{soln}
    This is essentially just our definition of Sylow $p$-subgroups:
    a Sylow $p$-subgroup is a $p$-subgroup of $G$ that is maximal under inclusion,
    and $H$, by definition of maximality,
    must be contained in a maximal (Sylow) $p$-subgroup of $G$.
\end{soln}

\begin{problem}
    Let $P$ and $Q$ be Sylow $p$-subgroups of a finite group $G$.
    Suppose the center $Z(P)$ of $P$ is contained in $Q$
    and is normal in $Q$. Prove that $Z(P) = Z(Q)$.
\end{problem}

\begin{soln}
    Let $H = N_G(Z(P))$ be the normalizer of $Z(P)$.
    Then $P$ and $Q$ are Sylow $p$-subgroups of $H$
    because the largest power of $p$ dividing $\abs H$ is at most
    the largest power of $p$ dividing $\abs G$,
    so there exists $h\in H$ such that $Q = hPh\inv$.
    At the same time, $hZ(P)h\inv = Z(P)$ because $h$ is in the normalizer of $Z(P)$.

    The conjugation-by-$h$ map $\phi$ from $P$ to $Q$ (i.e. $\phi(x) = hxh\inv$)
    is an isomorphism and it maps the center of $P$ bijectively to the center of $Q$:
    to verify this we see that $ab = ba$ for $a\in Z(P)$, $b\in P$
    implies $\phi(a)\phi(b) = \phi(b)\phi(a)$ where
    $\phi(a)\in hZ(P)h\inv$ and $\phi(b)\in Q$;
    conversely, $ab = ba$ for $a\in hZ(P)h\inv$, $b\in hPh\inv = Q$ implies
    $\phi\inv(a)\phi\inv(b) = \phi\inv(b)\phi\inv(a)$ where
    $\phi\inv(a)\in Z(P)$ and $\phi\inv(b)\in P$.
    In other words, $Z(Q) = hZ(P)h\inv = Z(P)$ as desired.
\end{soln}

\pagebreak

\begin{problem}
    Let $G$ be a group of order $p^2q$, where $p$ and $q$ are distinct primes.
    Prove that $G$ has a normal Sylow subgroup.
\end{problem}

\begin{soln}
    Let $n_p$ be the number of Sylow $p$-subgroups
    and $n_q$ be the number of Sylow $q$-subgroups, meaning
    \begin{align*}
        n_p &\equiv 1\bmod p,\quad n_p\mid q \\
        n_q &\equiv 1\bmod q,\quad n_q\mid p^2.
    \end{align*}
    If either $n_p$ or $n_q$ are $1$ then we are done,
    so assume otherwise. In this case $n_p = q$ and $n_q = p$ or $p^2$.
    Thus $p\mid q-1$ by the $1\bmod p$ constraint and
    $n_q\equiv 1\pmod q$ means that $q$ divides either
    $p-1$ or $p^2-1$, but $q\mid p^2-1$ encompasses both cases.
    That is, we want to find prime numbers $p$ and $q$ such that
    $p\mid q-1$ and $q\mid p^2-1$.
    Since $q$ divides $p^2-1 = (p-1)(p+1)$, it divides one of $p-1$ and $p+1$.
    However $q$ cannot divide $p-1$ because $p\mid q-1$ means that $q>p>p-1$,
    so we have $q\mid p+1$. This in turn implies $q\le p+1$,
    which, in addition to $p\le q-1$, implies that $p+1=q$.
    Thus the only pair of primes satisfying the divisibility constraints
    is $(p,q) = (2,3)$.

    Now, let $P_2$ and $P_3$ be respectively Sylow $2$- and $3$-subgroups of $G$.
    Their intersection $P_2\cap P_3$ must be the trivial group
    because its order must divide $\gcd(\abs{P_2},\abs{P_3}) = \gcd(4,3) = 1$.
    By the number-theoretic restrictions, $n_2$ is either $1$ or $3$
    and $n_3$ is either $1$ or $4$. If $n_3=1$ we are done, so suppose $n_3=4$.
    Any two distinct Sylow $3$-subgroups intersect only at the identity
    because Sylow $3$-subgroups are cyclic, so together
    the four Sylow $3$-subgroups contain $1 + 2\cdot 4 = 9$ elements of $G$.
    At this point there is only room for one Sylow $2$-subgroup,
    namely the identity element plus the three elements not in any of the
    Sylow $3$-subgroups, since any other hypothetical Sylow $2$-subgroup
    would intersect a Sylow $3$-subgroup at a non-identity element,
    which we argued above is impossible.
    Thus either $n_2$ or $n_3$ is $1$, which is what it means
    for either the Sylow $2$-subgroup or Sylow $3$-subgroup to be normal in $G$.
\end{soln}

\begin{problem}
    Let $G$ be a finite group, $P$ a Sylow $p$-subgroup of $G$.
    Suppose that $H$ is a subgroup of $G$ containing $N(P)$.
    Show that $(G:H)\quiv1\bmod p$.
\end{problem}

\begin{soln}
    We know that the normalizer of $P$ in $H$
    is the same as the normalizer of $P$ in $G$ because $H$ contains $N(P)$.
    Furthermore, since Sylow $p$-subgroups of $G$
    are all conjugate to one another, the number of Sylow $p$-subgroups
    is $(G:N(P))$, is $1\bmod p$.
    $P$ must also be a Sylow $p$-subgroup of $H$, so similarly
    $(H:N(P))\quiv1\bmod p$. Therefore
    \begin{equation*}
	(G:H) = (G:N(P))/(H:N(P))\quiv 1/1 = 1\bmod p.\qedhere
    \end{equation*}
\end{soln}

\begin{problem}
    Let $K$ be a normal subgroup of a finite group $G$.
    If $P$ is a Sylow $p$-subgroup of $K$, prove that $G = K\cdot N_G(P)$,
    where $N_G(P)$ is the normalizer of $P$ in $G$.
\end{problem}

\begin{soln}
    For every $g\in G$, the conjugate $gPg\inv$ of $P$
    must be a subgroup of $K$ because of the normality of $K$.
    Then $P$ and hence $gPg\inv$ are Sylow $p$-subgroups of $K$,
    so they are conjugate in $K$, i.e. there exists $k\in K$ such that
    $gPg\inv = kPk\inv$. Equivalently, $k\inv g\in N_G(P)$
    so $g\in kN_G(P)$, i.e. any $g\in G$ is contained in
    some coset of $N_G(P)$ by an element of $K$,
    which is what $K\cdot N_G(P) = G$ means.
\end{soln}

\begin{problem}
    Let $G$ be a finite group of order $80$.
    Show that $G$ is not simple.
\end{problem}

\begin{soln}
    Let $H\le G$ be a Sylow $2$-subgroup, so $\abs H = 16$.
    Then, $\abs G = 80$ does not divide $(G:H)! = 120$,
    so HW02 \#6(d) implies that $G$ is not simple.
\end{soln}

\pagebreak

\begin{problem}
    Suppose that the number $n_p$ of Sylow $p$-subgroups
    of a finite group $G$ satisfies $n_p\not\quiv1\bmod p^2$.
    Prove that there are two distinct Sylow $p$-subgroups $P$ and $Q$
    such that $\abs{P\cap Q} = \abs P/p$.
\end{problem}

\begin{soln}
    Fix an arbitrary Sylow $p$-subgroup $P$
    and let it act by conjugation on the set $X$ of Sylow $p$-subgroups.
    Then, the number of elementst of any orbit of this action
    is a power of $p$ (because it divides $\abs P$)
    This action has one fixed point (namely $P$)
    and $k$ orbits of size $p$; all other orbits have size
    $p^m$ for $m\ge 2$, so the sum of the sizes of the orbits is
    \begin{equation*}
	\abs X = n_p \quiv 1 + kp \bmod p^2,
    \end{equation*}
    which is not $1\bmod p^2$, so $p\nmid k$ and in particular $k\ne 0$.
    Pick $Q\in X$ whose orbit has size $p$.
    The stabilizer of $Q$ then has size $\abs P/p$ by orbit-stabilizer.
    If $pQp\inv = Q$ for some $p\in P$ then $p\in P\cap N_G(Q)$,
    but since the order of $p\in N_G(Q)$ is a power of $p$,
    by the lemma we proved in class, $p\in Q$. Therefore the stabilizer of $Q$
    is $P\cap Q$, so $\abs{P\cap Q} = \abs P/p$.
\end{soln}

\begin{problem}
    Let $G = N_0 \trianglerighteq N_1\trianglerighteq\cdots\trianglerighteq N_r=\{1\}$
    be a \vocab{composition series} for a finite group $G$, meaning that
    each $N_{i+1}$ is a normal subgroup of $N_i$, and that
    the factor groups $N_i/N_{i+1}$ are all simple.
    Suppose that one of the $N_i$s is a Sylow $p$-subgroup of $G$
    (for some prime $p$). Prove that $N_i$ must be normal in $G$.
\end{problem}

\begin{soln}
    Let $P = N_i$ be a Sylow $p$-subgroup of $G$;
    this means $P$ is a Sylow $p$-subgroup of $N_1$, \dots, $N_{i-1}$ as well.
    We prove that if $P$ is normal in $N_{i-j}$
    for some $1\le j\le i-1$, then $P$ is also normal in $N_{i-(j+1)}$
    via induction on $j$.

    For the base case $j=1$, we have that
    $P$ is normal in $N_{i-1}$ and $N_{i-1}$ is normal in $N_{i-2}$.
    Consider the conjugates of $P$ under $N_{i-2}$;
    these are the Sylow $p$-subgroups of $N_{i-2}$.
    The normality of $N_{i-1}$ means that conjugates of $P$
    stay subgroups of $N_{i-1}$, hence are Sylow $p$-subgroups;
    this means that $N_{i-2}$ also has $P$ as its only Sylow $p$-subgroup,
    hence $P$ is normal in $N_{i-2}$.

    The proof in the general case is similar;
    $P$ is normal in $N_{i-j}$, so conjugation by $N_{i-j-1}$
    sends $P$ to Sylow $p$-subgroups of $N_{i-j}$, which are still $P$,
    i.e. $P$ is fixed under conjugation by $N_{i-j-1}$ so it is normal in $N_{i-j-1}$.

    Thus taking for $j=i$ tells us that $P$ is normal in $N_0 = G$.
\end{soln}

\begin{problem}
    Suppose $G$ is a nontrivial finite $p$-group.
    \begin{enumerate}[(a)]
        \ii Show that the center of $G$ is nontrivial.
	\ii Show that $G$ is simple if and only if $\abs G = p$.
	\ii Let $\abs G = p^n$ and let
	$G = N_0\trianglerighteq N_1\trianglerighteq\cdots\trianglerighteq N_r
	= \{1\}$ be a composition series for $G$.
	Show that $r = n$ and $\abs{N_i} = p^{n-i}$ for $i=0,1,\dots,n$.
	\ii Show that we can find a composition series for $G$
	such that each $N_i$ is normal in $G$.
    \end{enumerate}
\end{problem}

\begin{soln}
    (a) The sizes of the conjugacy classes of $G$
    are divisors of $G$, hence powers of $p$, that sum to the order of $G$.
    The center is comprised of the fixed points under conjugation,
    the number of which has the same remainder mod $p$ as $\abs G$,
    namely zero. That is, $\abs{Z(G)}$ is a multiple of $p$
    and is at least $1$ because it contains the identity of $G$,
    so $Z(G)$ must be nontrivial.
    \\

    (b) If $\abs G > p$ then either $Z(G)\ne G$ and so $Z(G)$ is a nontrivial
    (not equal to $G$ or $\{1\}$) normal subgroup,
    or $Z(G) = G$ in which case $G$ is abelian; here any subgroup is normal
    and there exists a subgroup isomorphic to $\ZZ/p\ZZ$ by Cauchy's theorem
    (which is not equal to $G$ because its order is too small).

    On the other hand if $\abs G = p$ then any subgroup has order dividing
    $p$, hence the order of the subgroup is $1$ or $p$,
    which means that the subgroup is trivial or it is $G$.
    Since $G$ has no nontrivial proper subgroups it certainly can't have
    nontrivial proper \emph{normal} subgroups.
    \\

    (c) The orders of the $N_i$ are all powers of $p$
    so $N_{i-1}/N_i$ must have order $p$ in order to be simple
    due to part (b). This implies that $\abs{N_{i-1}} = \abs{N_i}/p$,
    so by induction we get $\abs{N_i} = p^{n-i}$ for all $i=0,1,\dots,n$
    and $r=n$ because the order reaches $1$ after $n$ divisions by $p$.
    \\

    (d) Let $\abs G = p^n$.
    We induct on $n$, with $n=1$ trivial because $\abs G = p$
    means $G$ is simple.

    In the general case, the center of $G$ is nontrivial
    and has order a power of $p$, so it has a subgroup $N$
    isomorphic to $\ZZ/p\ZZ$; this subgroup is normal in $G$
    because its elements commute with all the elements of $G$
    (due to being a subgroup of the center).
    Then $G/N = \ol G$ has order $p^{n-1}$, so it has a composition series
    of subgroups
    $\ol G\trianglerighteq\ol{N_1}\trianglerighteq\cdots\trianglerighteq\ol{N_r}=\{N\}$
    that are all normal in $G$ (where $\ol{N_r}$ consists of just the one coset of $N$)
    By the fourth isomorphism theorem, the subgroups $\ol N_i$ in this chain
    all correspond to normal subgroups $N_i$ of $G$ that contain $N$,
    and $N_{i-1}/N_i\cong \ol N_{i-1}/\ol{N_i}$ is simple for all $i=1,\dots,r$.
    Thus, lifting the $\ol{N_i}$ to normal subgroups of $G$,
    we get the composition series of $G$
    \begin{equation*}
        G = N_0\trianglerighteq N_1\trianglerighteq \cdots
	\trianglerighteq N_r = N \trianglerighteq \{1\}
    \end{equation*}
    with only normal subgroups of $G$.
\end{soln}










\end{document}
